% ch02-problem.tex – Chapter 2: problem statement, research question and goals (final round, Writer).
\ChapterMission{המשימה: להגדיר את הבעיה במספרים, לנסח שאלת מחקר והשערות ולקבוע יעדי תכן מדידים}
\ChapterIcon{warning}
\chapter{הגדרת הבעיה, שאלת המחקר והיעדים}\label{chap:problem}
\begin{scenebox}
צהריים. אותו מעבר, ועל המדרכה עומדת אישה אחת עם שקית קניות. האיש האדום דולק, והיא מעבירה את השקית מיד ליד.
מכונית חולפת, ואחריה שקט. היא מביטה ברמזור, אחר כך בכביש, ושוב ברמזור.
הרגל כבר מונחת על שפת המדרכה.
\end{scenebox}

ההתלבטות של הממתינה היחידה מגדירה את הבעיה: המתנה ארוכה מול כביש ריק, בלי איש שירסן את ההחלטה לחצות.
הפרק מכמת את הבעיה, מנסח את שאלת המחקר ואת ההשערות, וקובע יעדי תכן מדידים שמהם נגזרות הדרישות בפרק~\ref{chap:requirements}.

\section{הבעיה במספרים}
במעבר שנצפה בתל אביב חצו באדום \pct{13.5} ממי שהגיעו במופע האדום, וכל אדם נוסף שהמתין על המדרכה הקטין את הסיכויים לכך \cite{rosenbloom2009}; לפי החישוב מחדש, פי \num{0.63} לכל ממתין (סעיף~\ref{sec:critical}).
בשישה מעברים בהונג קונג חצו באדום \pct{18.5} מן המגיעים באדום; מי שהיה לו מלווה חצה פחות, וכל שנייה נוספת של אדום העלתה את סיכויי ההפרה \cite{zhu2021}.
גם זמן ההמתנה עצמו הוא גורם סיכון: הנטייה להפר עולה ככל שההמתנה מתארכת \cite{guo2011}, וכשההמתנה שונתה בניסוי שדה בין \unit{30}{s} ל-\unit{120}{s}, הציות ירד ככל שגדלה \cite{vanhouten2007}.
הבעיה היא אפוא הצירוף של הולך יחיד והמתנה ארוכה. איור~\ref{fig:t4} מסכם את המנגנון ומראה היכן המערכת המוצעת מתערבת: היא מקצרת את האדום להולך היחיד, ומחזקת בקמפיין את הנורמה של המתנה.

\begin{figure}[ht]
  \centering
  \inputfigure{figures/tikz/t4-social-loop}
  \caption{לולאת השליטה החברתית: ממתינים אחרים מרסנים, הולך יחיד ממתין לבדו, והמערכת מקצרת את האדום ומחזקת את הנורמה}
  \label{fig:t4}
\end{figure}

\section{שאלת המחקר וההשערות}
\begin{definitionbox}[שאלת המחקר]
האם קיצור האדום להולכי רגל ממתינים, בעדיפות להולך יחיד ובגבולות בטוחים, יחד עם ספירה לאחור במופע הפינוי וקמפיין שליטה חברתית, מפחית חצייה באדום בלי לפגוע בזרימת כלי הרכב?
\end{definitionbox}
\en{H1}: שיעור החצייה באדום בקרב המגיעים באדום יירד ביחס לנתוני הבסיס של אותו אתר.

\en{H2}: הירידה תהיה גדולה יותר בקרב הולכים יחידים, שאצלם השליטה החברתית חלשה.

\en{H3}: תוספת העיכוב לכלי הרכב תישאר בגבול שנקבע בדרישות.
ההשערות נבחנות בניסוי לפני ואחרי, ולכן חבילת העבודה האחרונה בפרק~\ref{chap:workplan} היא פיילוט עם מדידת בסיס.

\section{יעדי תכן מדידים}
טבלה~\ref{tab:problem} מציבה מול כל מדד את המצב הקיים ואת היעד. היעדים מנוסחים לפי עקרון \en{SMART}, כלומר ספציפיים, מדידים, ברי השגה, רלוונטיים ותחומים בזמן, ומועד היעד של כולם הוא סוף הפיילוט. ערכי היעד הם יעדי הצוות.

\begin{table}[ht]
  \centering
  \caption{הבעיה במספרים: מצב קיים לעומת יעדי התכן}
  \label{tab:problem}
  \small
  \begin{tabularx}{\linewidth}{lXX}
    \toprule
    \headrow \thead{יעד} & \thead{מצב קיים} & \thead{מצב רצוי} \\
    \midrule
    \en{GOAL-1} & \pct{13.5} מן המגיעים באדום חוצים באדום \cite{rosenbloom2009} & ירידה יחסית של \pct{25} לפחות מול נתוני הבסיס \\
    \en{GOAL-2} & זמן ההמתנה של הולך יחיד אינו נמדד \cite{rosenbloom2009} & נמדד בכל מופע, בדיוק שנקבע בדרישות \\
    \en{GOAL-3} & משך אדום קבוע, בלי קשר לממתינים & אדום מקוצר להולך יחיד, בלי אף הפרה של הגבולות הבטוחים \\
    \en{GOAL-4} & ספירה בזמן האדום אינה מותרת בתקן \cite{fhwa2025mutcd} & ספירה רק במופע הפינוי, בכל מופע \\
    \en{GOAL-5} & מצלמות עלולות לשמור תמונות של הולכים & אפס תמונות שמורות; ספירות וזמנים בלבד \\
    \en{GOAL-6} & עיכוב כלי הרכב בתזמון הקבוע & תוספת עיכוב ממוצעת של \pct{5} לכל היותר \\
    \bottomrule
  \end{tabularx}
\end{table}

איור~\ref{fig:t1} ממקם את היעדים במרחב: אזור ההמתנה, אזור החישה בשפת המדרכה, ראשי האות עם הספירה במופע הפינוי, הבקר ותצוגת הקמפיין.

\begin{figure}[ht]
  \centering
  \inputfigure{figures/tikz/t1-crosswalk-zones}
  \caption{מעבר החצייה החכם במבט על: אזור ההמתנה, אזור החישה, ראשי אות עם ספירה במופע הפינוי, בקר האדום המסתגל ותצוגת הקמפיין}
  \label{fig:t1}
\end{figure}
